%%%%%%%%%%%%%%%%%%%%%%%%%%%%%%%%%%%%%%%%%%%%%%%%%%%%%%%%%%%%%%%%%%%%%%%%
%%%%%%%%%%%%%%%%%%%%%% Simple LaTeX CV Template %%%%%%%%%%%%%%%%%%%%%%%%
%%%%%%%%%%%%%%%%%%%%%%%%%%%%%%%%%%%%%%%%%%%%%%%%%%%%%%%%%%%%%%%%%%%%%%%%

%%%%%%%%%%%%%%%%%%%%%%%%%%%%%%%%%%%%%%%%%%%%%%%%%%%%%%%%%%%%%%%%%%%%%%%%
%% NOTE: If you find that it says                                     %%
%%                                                                    %%
%%                           1 of ??                                  %%
%%                                                                    %%
%% at the bottom of your first page, this means that the AUX file     %%
%% was not available when you ran LaTeX on this source. Simply RERUN  %%
%% LaTeX to get the ``??'' replaced with the number of the last page  %%
%% of the document. The AUX file will be generated on the first run   %%
%% of LaTeX and used on the second run to fill in all of the          %%
%% references.                                                        %%
%%%%%%%%%%%%%%%%%%%%%%%%%%%%%%%%%%%%%%%%%%%%%%%%%%%%%%%%%%%%%%%%%%%%%%%%

%%%%%%%%%%%%%%%%%%%%%%%%%%%% Document Setup %%%%%%%%%%%%%%%%%%%%%%%%%%%%

% Don't like 10pt? Try 11pt or 12pt
\documentclass[10pt]{article}
\usepackage{times}
\usepackage[T1]{fontenc}
\usepackage{lmodern}
\renewcommand*\familydefault{\sfdefault} %% Only if the base font of the document is to be sans serif

% This is a helpful package that puts math inside length specifications
\usepackage{%
	calc,%
	%enumitem,%
	hyperref%
}

%% This gives us fun enumeration environments. compactenum will be nice.
\usepackage{paralist}
% Change the margin here to 
\setdefaultleftmargin{1.5em}{}{}{}{}{}

% Layout: Puts the section titles on left side of page
\reversemarginpar

%
%         PAPER SIZE, PAGE NUMBER, AND DOCUMENT LAYOUT NOTES:
%
% The next \usepackage line changes the layout for CV style section
% headings as marginal notes. It also sets up the paper size as either
% letter or A4. By default, letter was used. If A4 paper is desired,
% comment out the letterpaper lines and uncomment the a4paper lines.
%
% As you can see, the margin widths and section title widths can be
% easily adjusted.
%
% ALSO: Notice that the includefoot option can be commented OUT in order
% to put the PAGE NUMBER *IN* the bottom margin. This will make the
% effective text area larger.
%
% IF YOU WISH TO REMOVE THE ``of LASTPAGE'' next to each page number,
% see the note about the +LP and -LP lines below. Comment out the +LP
% and uncomment the -LP.
%
% IF YOU WISH TO REMOVE PAGE NUMBERS, be sure that the includefoot line
% is uncommented and ALSO uncomment the \pagestyle{empty} a few lines
% below.
%

%% Use these lines for letter-sized paper
%\usepackage[paper=letterpaper,
            %%includefoot, % Uncomment to put page number above margin
            %marginparwidth=1.2in,     % Length of section titles
            %marginparsep=.05in,       % Space between titles and text
            %margin=1in,               % 1 inch margins
            %includemp]{geometry}

%% Use these lines for A4-sized paper
%\usepackage[paper=a4paper,
%            includefoot, % Uncomment to put page number above margin
%            marginparwidth=30.5mm,    % Length of section titles
%            marginparsep=1.5mm,       % Space between titles and text
%            margin=25mm,              % 25mm margins
%            includemp]{geometry}

%% Use these lines for maximum space utilization for CV or resume
\usepackage[paper=a4paper,
            includefoot, % Uncomment to put page number above margin
            marginparwidth=27.5mm,    % Length of section titles
            marginparsep=0mm,       % Space between titles and text
            margin=16mm,              % 16mm margins
            includemp]{geometry}

%% More layout: Get rid of indenting throughout entire document
\setlength{\parindent}{0in}

%% Reference the last page in the page number
%
% NOTE: comment the +LP line and uncomment the -LP line to have page
%       numbers without the ``of ##'' last page reference)
%
% NOTE: uncomment the \pagestyle{empty} line to get rid of all page
%       numbers (make sure includefoot is commented out above)
\usepackage{fancyhdr,lastpage}

\pagestyle{fancy}
\pagestyle{empty}      % Uncomment this to get rid of page numbers
\fancyhf{}\renewcommand{\headrulewidth}{0pt}
\fancyfootoffset{\marginparsep+\marginparwidth}
\newlength{\footpageshift}
\setlength{\footpageshift}
          {0.5\textwidth+0.5\marginparsep+0.5\marginparwidth-2in}
\lfoot{\hspace{\footpageshift}%
       \parbox{4in}{\, \hfill %
                    \arabic{page} of \protect\pageref*{LastPage} % +LP
%                    \arabic{page}                               % -LP
                    \hfill \,}}

% Finally, give us PDF bookmarks
\usepackage{color,hyperref}
\definecolor{darkblue}{rgb}{0.0,0.0,0.3}
\hypersetup{colorlinks,breaklinks,
            linkcolor=darkblue,urlcolor=darkblue,
            anchorcolor=darkblue,citecolor=darkblue}

%%%%%%%%%%%%%%%%%%%%%%%% End Document Setup %%%%%%%%%%%%%%%%%%%%%%%%%%%%


%%%%%%%%%%%%%%%%%%%%%%%%%%% Helper Commands %%%%%%%%%%%%%%%%%%%%%%%%%%%%

% The title (name) with a horizontal rule under it
%
% Usage: \makeheading{name}
%
% Place at top of document. It should be the first thing.
\newcommand{\makeheading}[1]%
        {\hspace*{-\marginparsep minus \marginparwidth}%
         \begin{minipage}[t]{\textwidth+\marginparwidth+\marginparsep}%
                {\large \bfseries #1}\\[-0.15\baselineskip]%
                 \rule{\columnwidth}{1pt}%
         \end{minipage}}

% The section headings
%
% Usage: \section{section name}
%
% Follow this section IMMEDIATELY with the first line of the section
% text. Do not put whitespace in between. That is, do this:
%
%       \section{My Information}
%       Here is my information.
%
% and NOT this:
%
%       \section{My Information}
%
%       Here is my information.
%
% Otherwise the top of the section header will not line up with the top
% of the section. Of course, using a single comment character (%) on
% empty lines allows for the function of the first example with the
% readability of the second example.
\renewcommand{\section}[2]%
        {\pagebreak[2]\vspace{0.7\baselineskip}%
         \phantomsection\addcontentsline{toc}{section}{#1}%
         \hspace{0in}%
         \marginpar{
         \raggedright \bfseries #1}#2}

% An itemize-style list with lots of space between items
\newenvironment{outerlist}[1][\enskip\textbullet]%
        {\begin{enumerate}[#1]}{\end{enumerate}%
         \vspace{-1\baselineskip}}

% An environment IDENTICAL to outerlist that has better pre-list spacing
% when used as the first thing in a \section
\newenvironment{lonelist}[1][\enskip\textbullet]%
        {\vspace{-\baselineskip}\begin{list}{#1}{%
        \setlength{\partopsep}{0pt}%
        \setlength{\topsep}{0pt}}}
        {\end{list}\vspace{-.6\baselineskip}}

% An itemize-style list with little space between items
\newenvironment{innerlist}[1][\enskip\textbullet]%
        {\begin{compactenum}[#1]}{\end{compactenum}}

% To add some paragraph space between lines.
% This also tells LaTeX to preferably break a page on one of these gaps
% if there is a needed pagebreak nearby.
\newcommand{\blankline}{\quad\pagebreak[2]}

\usepackage{fontawesome}

%%%%%%%%%%%%%%%%%%%%%%%% End Helper Commands %%%%%%%%%%%%%%%%%%%%%%%%%%%

%%%%%%%%%%%%%%%%%%%%%%%%% Begin CV Document %%%%%%%%%%%%%%%%%%%%%%%%%%%%

\begin{document}

\makeheading{\huge \textsc{Bidhov Bizar} \\
\rule{\textwidth}{1pt}
\large Software Engineer | Cisco | IISc | RIT}

\section{Contact Information}
% NOTE: Mind where the & separators and \\ breaks are in the following
%       table.
%
% ALSO: \rcollength is the width of the right column of the table
%       (adjust it to your liking; default is 1.85in).
\newlength{\rcollength}\setlength{\rcollength}{2.5in}%
%
\begin{tabular}[t]{@{}p{\textwidth-\rcollength}p{\rcollength}}
Lotus Krest Phase 1       & +91 7012-556-335 \\
Green Domain Layout, Brookefield   &                    bbizar@cisco.com   \\
Bangalore, KA 560037, India          & \href{https://www.linkedin.com/in/bidhovbizar/}{https://www.linkedin.com/in/bidhovbizar}\\
\end{tabular}

\section{Professional Summary}
Software Automation \& Systems Engineer with expertise in distributed systems, networking, and cloud-managed infrastructure within Cisco's Data Center BU.
Experienced in building and scaling automation scripts and testing infrastructure and networked systems, with a strong focus on reliability and scalability.
Proven ability to debug end-to-end system behavior across hardware, firmware, and microservices-based architectures using Graphana, APIs, and trace-based debugging.
Skilled in improving CI/CD pipelines, owning features end-to-end, and driving reliability improvements across large-scale validation and execution environments.
Proficient in Go-based backend services and production codebases.

\section{Interests}
Software Development and Automation, Distributed Systems and Networking, L2/L3 Network Protocols, Datacenter Life cycle management, Edge Intelligence, Compute host and BMC, Jenkins and python automation, Microservice and Kubernetes, Generative AI

\section{Skills}
\begin{tabular}{l l}
\textbf{Programming Languages:} & Go, Python, C, C++, Bash, Groovy\\
\textbf{Domains:} &UCS Servers, Nexus Switch, Distributed Systems, Platform Engineering \\
\textbf{Tools \& Technologies:} & GNS3, Docker, Kibana, Graphana, WireShark, Mininet, TRex, IXIA\\
\textbf{Operating System:} & {\sc Linux}, Windows, CentOS, RHEL, SONiC, NX-OS\\
\textbf{AI Tools:} & Copilot, Cursor, Codex, Chatgpt, Avante\\
\end{tabular}

\section{Professional Experience}
\textbf{Automation Engineer/Systems Engineer, Intersight BU, Cisco} \hfill  Aug 2023 -- Present

\quad \textbf{Automation Engineering \& Platform / Infrastructure  }
\begin{compactitem}
\item Designed and implemented 15+ \textbf{test plans} and scalable automation frameworks, including \textbf{emulator-based execution}, enabling validation across hardware and multiple cloud environments.
\item Developed and maintained \textbf{end-to-end automation workflows} for firmware upgrades, OS validation, server policies, adapter configuration, and platform health checks, integrated with \textbf{Jenkins-based CI/CD pipelines}.
\item Developed \textbf{test infrastructure workflows} including  testbed orchestration, and environment provisioning, improving reliability and repeatability of large-scale validation setups.
\item Diagnosed and resolved 40+ sev2 complex \textbf{platform and system issues} spanning firmware, inventory, discovery/claim flows, alarm schemas, and deployment semantics across \textbf{ISM, IMM, FI, appliance, and cloud environments}.
\item Enhanced script to improve\textbf{observability and debugging} by improving alarm verification, tech-support collection, and RCA tooling, significantly reducing time-to-diagnose for failures.
\end{compactitem}

\quad \textbf{Systems, Backend \& Debugging}
\begin{compactitem}
\item Worked across \textbf{automation infra, platform, and Intersight backend}, tracing issues into developer-owned services and contributing fixes, including exposure to \textbf{Go-based systems}.
\item Performed deep \textbf{root cause analysis} across distributed systems, blocking critical defects from progressing into staging and production validation cycles.
\item Led bring-up, recovery, and operationalization of \textbf{lab and hardware environments} (UCSC/UCSX M7/M8, Fabric Interconnects, GEL hosts 7.9/8.10), ensuring platform readiness for qualification.
\end{compactitem}

\quad \textbf{CI/CD, Framework Stability \& Release Engineering}
\begin{compactitem}
\item Led \textbf{CI/CD and build system improvements}, including resolving static analysis failures, upgrading \textbf{Python runtimes (3.7 $\rightarrow$ 3.12)}, fixing Docker dependencies, and stabilizing Jenkins pipelines.
\item Identified and fixed \textbf{framework-level issues} (pylint recursion limits, schema mismatches, serialization gaps), improving robustness of large-scale automation and analysis workflows.
\item Drove \textbf{release and patch qualification} across multiple platform variants (Compute, FI infra bundles 4.2.x to 6.0.2.x), ensuring readiness for staging and production releases.
\item Improved execution reliability by resolving issues in \textbf{pipeline orchestration, logging, report portal integration, and execution host compatibility}.
\end{compactitem}

\quad \textbf{Developer Productivity, Tooling \& Leadership}
\begin{compactitem}
\item Built internal \textbf{AI-assisted tools} including a test report summarizer, daily sanity bot, and integrations with \textbf{CDets and Kibana} to accelerate failure analysis.
\item Developed a \textbf{Jira-Webex integration tool} for automated reporting and tracking, improving visibility and reducing coordination overhead across teams.
\item Contributed to \textbf{AI/DeepInsight initiatives} (Bugscribe, test plan generation), presenting solutions to internal teams and leadership.
\item Mentored new engineers through \textbf{knowledge transfer sessions}, onboarding, and technical workshops on coding practices and tooling.
\end{compactitem}

\textbf{Software Engineer, Data Center BU, Cisco} \hfill July 2022 -- July 2023

\quad \textbf{L3 Team}
\begin{compactitem}
\item Designed and implemented \textbf{CLI-based features} based on customer requirements across multiple release cycles.
\item Debugged complex \textbf{L2/L3 networking issues} (ISIS, BFD, VRF, ARP) and forwarding behavior in large-scale deployments.
\item Developed expertise in \textbf{control plane and inter-process communication} between supervisor and line cards.
\item Led \textbf{root cause analysis} of customer-reported issues and collaborated cross-functionally to resolve \textbf{release-blocking defects}.
\item Set up and validated \textbf{Nexus switch topologies (fat-tree)} for testing and debugging.
\item Conducted internal sessions on \textbf{CLI architecture and DME integration}, and contributed to technical documentation.
\end{compactitem}

\quad \textbf{L2 Team} \hfill Jan 2022 -- July 2022
\begin{compactitem}
\item Developed and automated \textbf{SPAN/ERSPAN validation frameworks} for Cisco ASICs in virtualized environments.
\item Built and debugged \textbf{containerized (Docker) network testing workflows}.
\item Investigated issues in \textbf{L2 forwarding and traffic mirroring}, and performed detailed \textbf{packet walk analysis} on N7k/N9k platforms.
\end{compactitem}

%\textbf{Software Engineering Intern, Data Center BU, Cisco} \hfill May 2021 -- July 2021
%\begin{compactitem}
%\item Analyzed and debugged \textbf{convergence delays} during switch failover (standby to active transition).
%\item Identified performance bottlenecks affecting \textbf{control plane convergence} in L2/L3 components.
%\item Developed Python-based tools to \textbf{automate detection of convergence issues} and improve debugging efficiency.
%\item Delivered internal training sessions on automation techniques for network performance analysis.
%\end{compactitem}

\vspace{-6pt}
\section{Projects Experience}
\vspace{-.6cm}
\begin{outerlist}
\item Built an \textbf{AI-based test report summarizer} integrating logs, metrics, and scribe the cdet defects.
Assisted in the development AI tools for \textbf{test plan generation, defect analysis, and traceid analysis}.
\item {\bf Networking:}\\
Constructed network using {\bf Mininet} with multiple {\bf TRex Traffic Generator} at hosts running stateless, stateful and Client Server L7 traffic and obtained telemetry.[2021]\\
Constructed {\bf \href{https://github.com/bidhovbizar/DCN-Simulation/tree/master/PythonCode}{Simulation Framework}} in {\bf Python} for large scale Data Center in Fat-Tree Topology. [2020]  \\ 
Emulated Fat-Tree topology in Mininet to collect the telemetry data and observe congestion using {\bf Ryu} Controller and {\bf OpenVSwitches}.[2020]
\vspace{-6pt}
\end{outerlist}


\section{Education}
\textbf{M.~Tech.(Research), Indian Institute of Science}% 
\hfill 07/2018-Present\\ %\vspace{-.2cm}
Electrical Communication Engineering, Network Lab
        \begin{innerlist}
        \item[$\diamond$] Dissertation: \emph{``Empirical study of load generation for Data Center Network''}
       % \item[$\diamond$] Advisor:     {\href{https://ece.iisc.ac.in/~parimal/}{Parimal Parag}}
        
        \end{innerlist}
\vspace{2pt}

\textbf{B.~Tech., Rajiv Gandhi I.T., Kottayam}%
	\hfill 07/2012-07/2016\\ %\vspace{-.2cm}
Electrical Engineering\\
Class Rank : 7/72
\vspace{2pt}

\textbf{Higher Secondary and High School, Kendriya Vidyalaya, KGQ}%
	\hfill 06/2010-07/2012\\ %\vspace{-.2cm}
Science Stream \\
Percentage : 94 and 92 respectively
\vspace{2pt}

\section{Relevant Courses}
\textbf{Electrical Engineering}: 
 Digital Communications, Communication Networks, Real Time System\\
\textbf{Mathematics}: 
Probability, Convex Optimization, Queuing Thoery.\\
\textbf{Data Science}: 
 Machine Learning, Practical Data Science, Pattern Recognition \& Neural Network

\section{Honors \& Awards}
\vspace{-.5cm}
\begin{compactitem}
\item All India rank 1276 in ECE stream in GATE, 2018.
%\item Project shortlisted for AIYEHUM,IEEE R10.[2016]
%\item Shortlisted for Mathematics camp For INMO (2011,2012)
\item $99^{th}$ rank in {Junior Mathematics Olympiad}, 2010.
\end{compactitem}
\renewcommand{\labelenumi}{\arabic{enumi}.}

%\section{Memberships}
%IEEE, Computer Society, WIE
%\blankline

\section{References}\vspace{-5mm}
\begin{center}
\begin{tabular}{ll}
\textbf{Anirudh Sripuram}  \\
Leader, Software Engineering\\
Intersight Compute BU\\
Cisco Systems India Pvt Ltd, Bengaluru\\
E-mail: asripura@cisco.com\\[1cm]
\end{tabular}
\begin{tabular}{ll}
\textbf{Prasanth Rajan Sivakumar} \\
Leader, Software Engineering\\
Intersight Compute BU\\
Cisco Systems India Pvt Ltd, Bengaluru\\
E-mail: psivakum@cisco.com\\[1cm]
\end{tabular}

\end{center}

\end{document}

%%%%%%%%%%%%%%%%%%%%%%%%%% End CV Document %%%%%%%%%%%%%%%%%%%%%%%%%%%%%
